\section{DLT kao usmereni aciklični graf}

U dosadašnjem delu rada definisane su osnovne karakteristike usmerenog acikličnog grafa i motivacija za korišćenje takve strukture u distribuiranim sistemima. U ovom poglavlju fokus će biti na primeni DAG-a kao organizacione strukture u distribuiranim glavnim knjigama. Biće diskutovano kako DAG utiče na arhitekturu, konsenzus i ponašanje DLT sistema.
%\colorbox{yellow}{DODATI:~U ODNOSU NA BLOKČEJN}. 

Kod blokčejn baziranih DLT sistema transakcije se grupišu u blokove koji se zatim dodaju u linearnu, sekvencijalnu listu blokova, pri čemu svaki novi blok pokazuje na prethodni kroz kriptografski heš, formirajući neprekinuti lanac podataka. Pre dodavanja novog bloka u lanac, njegova ispravnost se mora potvrditi putem konsenzus mehanizma. Kod protokola kao što su dokaz posla (engl. \textit{Proof-of-Work}) i dokaz uloga (engl. \textit{Proof-of-Stake}), neophodna je saglasnot većine čvorova mreže za validaciju bloka. Ovakva struktura dovodi do ograničenja u pogledu brzine obrade transakcija, jer se novi blok može dodati samo nakon što se prethodni potrvdi~\cite{university mitosis blockchain vs blockDAG}.

Glavni ciljevi korišćenja DAG-a u ovakvim sistema jesu poboljšanje skalabilnosti, brzine i troškova transakcija tradicionalnih blokčejn arhitektura~\cite{blockchain council role of DAG}.

\subsection{Struktura DLT-a kao usmerenog acikličnog grafa}

DLT sistemi bazirani na usmerenom acikličnom grafu razlikuju tri tipa elemenata u grafu: transakcije, događaje i blokove. Transakcija predstavlja osnovnu jedinicu koja menja stanje sistema, događaj služi za razmenu informacija o redosledu i komunikaciji između čvorova, dok blok može objediniti više transakcija~\cite{DAG in Smart Mobility}.

Na osnovu izbora osnovne jedinice reprezentacije, DAG DLT sistemi se dalje dele na \textit{blockless} i \textit{block-based} arhitekture. U prvim sistemima elementi u grafu su pojedinačne transakcije ili događaji koji se neposredno rasprostiru kroz graf, dok u drugim sistemima graf čine blokovi koji zahtevaju dodatno procesiranje od strane određenih čvorova pre dodavanja u graf~\cite{DAG in Smart Mobility}.

Kod DAG DLT sistema svaka transakcija postaje deo grafa i povezuje se sa jednom ili više prethodnih transakcija kako bi se potvrdila njihova validnost i istorija. Ovakav pristup omogućava da se istovremeno više transakcija izvršava i potvrđuje paralelno, rezultujući većim protokom transakcija i znatno kraćim vremenom obrade podataka. Takođe, ovo smanjuje potrebu za posebnim čvorovima koji validiraju transakcije, takozvanih \enquote{rudara}, i omogućava da sama aktivnost učesnika mreže doprinosi grafu. Međutim, u nekim sistemima i dalje postoje specifični čvorovi koji učestvuju u konsenzusu i finalizaciji transakcija~\cite{university mitosis blockchain vs blockDAG}.

%Kod DAG DLT sistema svaka transakcija postaje deo grafa i povezuje se sa jednom ili više prethodnih transakcija kako bi se potvrdila njihova validnost i istorija. Ovakav pristup omogućava da se istovremeno više transakcija izvršava i potvrđuje paralelno, rezultujući većim protokom transakcija i znatno kraćim vremenom obrade podataka~\cite{university mitosis blockchain vs blockDAG}.

%DAG bazirani DLT sistemi doprinose tome da validacija transakcija bude distribuiranija i skalabilnija u odnosu na klasične blokčejn sisteme. Nasuprot blokčejnu, u DAG strukturi, svaka transakcija pri svom unošenju u graf može potvrditi druge transakcije, čime sistem smanjuje potrebu za posebnim čvorovima koji validiraju transakcije, takozvanih \enquote{rudara}, i omogućava da sama aktivnost učesnika mreže doprinosi grafu.

Takođe, kod \textit{Proof-of-Work} blokčejn sistema postoji mogućnost nastanka takozvanog siročad čvora (engl. \textit{orphan block}) koji predstavlja blok koji ne pripada najdužem lancu blokova. Kada više čvorova rudara istovremeno validiraju i pošalju novi blok u mrežu, samo jedan blok je moguće priključiti najdužem lancu, dok će ostali, siročad čvorovi, biti odbačeni, rezultujući viškom potrošnje računarske snage~\cite{investopedia orphan block}. Sa druge strane, kod DAG baziranih DLT sistema ne postoje odbačeni čvorovi usled ovakvih situacija, što doprinosi manjim troškovima transakcija.

Ovakav način organizacije distribuirane glavne knjige omogućava efikasnije korišćenje resursa mreže, brži protok podataka i potencijalno niže troškove transakcija, ali takođe postavlja specifične zahteve za konsenzus algoritme i sigurnosne protokole koji moraju da održe funkcionisanje takve mreže.

\subsection{Konsenzus u DAG baziranim DLT sistemima}

U distribuiranim glavnim knjigama konsenzus predstavlja mehanizam kojim se usaglašava zajedničko stanje sistema među decentralizovanim čvorovima, to jest, način na koji se postiže dogovor o validnim transakcijama i njihovom redosledu. U tradicionalnim blokčejn sistemima, kosenzus se oslanjao na sekvencijalno dodavanje blokova uz pomoć protokola kao što su \textit{Proof-of-Work} ili \textit{Proof-of-Stake}, ali ove metode nisu pogodne za podatke koji nisu strukturirani kao linearan lanac blokova. DAG bazirani DLT sistemi zahtevaju konsenzus protokole koji omogućuju paralelno vršenje i dodavanje transakcija u mrežu~\cite{blockeden DAG}. %U nastavku će biti prikazano nekoliko pristupa rešavanju ovih problema.

Osnovni cilj konsenzus mehanizama u DAG baziranim sistemima jeste da obezbede jedinstven i konzistentan pogled na istoriju transakcija, kao i da spreče konflikte poput dvostruke potrošnje (engl. \textit{double spending}), uprkos viosokom stepenu paralelizma. Umesto sekvencijalnog potvrđivanja transakcija, konsenzus se u ovim sistemima često postiže akumulacijom potvrda ili glasova kroz samo strukturu grafa, gde nove transakcije potvrđuju prethodne~\cite{DAG in Smart Mobility}.

Jedan od ključnih principa kod mnogih DAG baziranih DLT sistema jeste da učesnici mreže, prilikom dodavanja nove transakcije u graf, preuzimaju deo odgovornosti za validaciju postojećih transakcija. Na taj način, proces validacije postaje distribuiran između svih aktivnih čvorova, a ne centralizovan oko nekoliko rudara ili validatora, kao što je to slučaj kod \textit{Proof-of-Work} ili \textit{Proof-of-Stake} sistema.

Konsenzus u DAG sistemima se može ostvariti na različite načine, u zavisnosti od konkretne implementacije. Kod nekih sistema, čvorovi probabilistički potvrđuju transakcije, a njihova verovatnoća konačne potvrde raste sa brojem potvrda. Drugi sistemi koriste mehanizme glasanja, gde određeni čvorovi imaju pravo da učestvuju u donošenju odluka o validnosti i redosledu transakcija, pri čemu težina njihovog glasa može zavisiti od različitih kriterijuma, kao što je količina uloženih sredstava ili reputacija u mreži.

Konsenzus mehanizmi kod DAG baziranih DLT sistema često teže eliminaciji energetskih troškova karakterističnih za \textit{Proof-of-Work} protokole, kao i smanjenju latencije u potvrđivanju transakcija. Međutim, ovakav pristup uvodi nove izazove, posebno u pogledu bezbednosti mreže, otpornosti na zlonamerne učesnike, takozvane vizantijske otkaze, i potrebe za dodatnim mehanizmima koji sprečavaju manipulaciju grafom u uslovima niskog mrežnog opterećenja~\cite{DAG in Smart Mobility}.

Zbog navedenih razlika, konsenzus u DAG baziranim DLT sistemima predstavlja aktivno istraživačko polje, sa različitim predloženim modelima i implementacijama. Konkretni primeri ovih mehanizama i sistema koji ih koriste biće detaljnije razmotreni u narednom poglavlju rada.